\documentclass[10pt,a4paper]{amsart}
\usepackage[margin=19mm]{geometry}
\usepackage{amsmath,amssymb,mathtools}
\usepackage[scr=rsfs,cal=euler]{mathalfa}
\usepackage{xfrac}
\usepackage[unicode,hidelinks,bookmarks=false]{hyperref}
\newcommand{\ie}{i.e.\ }
\newcommand{\cf}{cf.\ }
\newcommand{\resp}{respectively\ }
\newcommand{\Subsection}{\subsection}
\providecommand{\nobreakdash}{\nobreak\hbox{-}}
\setlength{\emergencystretch}{2em}
\multlinegap0pt
\usepackage{bm}
\usepackage{bbm}
\usepackage{dsfont}
\usepackage{xspace}
\usepackage{cancel}
\usepackage{mathtools}
\usepackage{amsthm}
\makeatletter
\@ifundefined{theorem}{\newtheorem{theorem}{Theorem}}{}
\@ifundefined{lemma}{\newtheorem{lemma}[theorem]{Lemma}}{}
\@ifundefined{proposition}{\newtheorem{proposition}[theorem]{Proposition}}{}
\makeatother
\newcommand{\K}{\mathbb K}
\newcommand{\algA}{\mathcal A}
\newcommand{\ev}{\operatorname{ev}}
\newcommand{\id}{\operatorname{Id}}
\title[Fixed-Lowering $\mathfrak{sl}\_2$-Triples for Laurent-Shift O]{Fixed-Lowering $\mathfrak{sl}\_2$-Triples for Laurent-Shift Operators: Exact Stencil Endpoints and Recurrence Locality}
\author{Kyle Singh}
\date{Source: arXiv:2608.09962v1 (CC BY 4.0)}
\begin{document}
\maketitle

\noindent\emph{Source and licence.} Fixed-Lowering $\mathfrak{sl}_2$-Triples for Laurent-Shift Operators: Exact Stencil Endpoints and Recurrence Locality. Version arXiv:2608.09962v1 (CC BY 4.0), distributed under the Creative Commons Attribution 4.0 International licence, as recorded in the arXiv licence metadata for this exact version and shown on the abstract page https://arxiv.org/abs/2608.09962v1. Full rights notice: licenses/arxiv-2608.09962-CC-BY-4.0.txt.\par\smallskip

\noindent\textit{Editorial scope.} Counted unit: the Theorem whose source label is \texttt{thm:classification} in the article (source0.tex lines 250--278, source label \texttt{thm:classification}), delivered together with the article's own complete original proof of it (source0.tex lines 280--384, one \texttt{proof} environment of 105 lines). No mathematics is written, repaired or re-derived: statement and proof are the article's own text line for line. Required context retained uncounted: prop:ore (lines 171--182); eq:sl2-relations (lines 228--234); lem:evaluation (lines 236--244). Every definition, display and label of those lines is retained unchanged. Omitted from this package and not redistributed: 1--170, 183--227, 235--235, 245--249, 279--279, 385--1394. Nothing inside a retained excerpt is omitted or altered: every retained body is delivered exactly as the article's lines are written, so no omission receipt of any kind accompanies this package. Citations: the retained excerpts carry no citation command at all, so no bibliography entry of the article is transcribed and no reference is redistributed. Reference adaptation: none was needed and none was made; every reference inside the delivered unit resolves inside it, so no unresolved reference or citation is produced. Printed slips kept literal: no printed word, symbol, hypothesis or number of the counted statement or of its proof was altered. Layout and class: the article's own class is \texttt{amsart} with packages including geometry, mathtools, amssymb, booktabs, tabularx, array, enumitem, microtype, hyperref; the wrapper is the lane's pinned \texttt{amsart} editorial wrapper at 10pt on A4, which restates the theorem-like environments the retained text needs and adds the article's own macro definitions that the retained text needs (\texttt{\textbackslash K}, \texttt{\textbackslash algA}, \texttt{\textbackslash ev}, \texttt{\textbackslash id}), with extra wrapper packages (bm, bbm, dsfont, xspace, cancel, mathtools). The delivered numbering is the unit's own. Assets: no retained span contains a figure environment, a graphics-inclusion command, a PDF-page inclusion, a table or a TikZ picture; no image file is copied into this package, the package declares no asset dependency and no image file is redistributed with it. Licence: version arXiv:2608.09962v1 of this article is distributed under the Creative Commons Attribution 4.0 International licence, as recorded in the arXiv licence metadata for this exact version and shown on the abstract page https://arxiv.org/abs/2608.09962v1. A full rights notice is delivered at licenses/arxiv-2608.09962-CC-BY-4.0.txt. Characters: every retained excerpt is pure ASCII, so the delivered engine, XeLaTeX with its default Latin Modern text font, reports no missing character.
\begin{proposition}\label{prop:ore}
The operators $D$ and $X$ satisfy $[D,X]=1$.  Moreover, $\algA$ is isomorphic to the Ore extension $R[X;-\partial_D]$, in which
\[
 Xr=rX-r',\qquad r\in R.
\]
Consequently every $T\in\algA$ has a unique normal form
\[
 T=\sum_{j=0}^{N}X^j r_j(D),
 \qquad
 r_j(D)\in R.
\]
\end{proposition}

\begin{equation}\label{eq:sl2-relations}
 [H,E]=2E,
 \qquad
 [H,F]=-2F,
 \qquad
 [E,F]=H.
\end{equation}

\begin{lemma}\label{lem:evaluation}
The evaluation homomorphism
\[
 \ev_0:R=\K[D,(D+1)^{-1}]\longrightarrow\K,
 \qquad
 r\longmapsto r(0),
\]
has kernel $DR$.
\end{lemma}

\begin{theorem}\label{thm:classification}
Let $E,H\in\algA$ and fix $F=-D=1-S$.  Then $(E,H,F)$ is an $\mathfrak{sl}_2$-triple if and only if there exist unique $\lambda\in\K$ and $g(D)\in R$ such that, with
\[
 Y=X+g(D),
\]
one has
\begin{equation}\label{eq:normal-form}
 H=2YD-\lambda,
 \qquad
 E=Y^2D-\lambda Y.
\end{equation}
For every $g\in R$, the assignment
\begin{equation}\label{eq:shear-automorphism}
 \sigma_g(D)=D,
 \qquad
 \sigma_g(X)=X+g(D)
\end{equation}
defines a $\K$-algebra automorphism of $\algA$ with inverse $\sigma_{-g}$.  If
\[
 H^{(0)}_\lambda=2XD-\lambda,
 \qquad
 E^{(0)}_\lambda=X^2D-\lambda X,
\]
then every fixed-lowering triple is uniquely
\[
 (E,H,F)=\bigl(\sigma_g(E^{(0)}_\lambda),
 \sigma_g(H^{(0)}_\lambda),-D\bigr).
\]
\end{theorem}

\begin{proof}
Suppose $(E,H,F)$ satisfies \eqref{eq:sl2-relations}.  Write
\[
 H=\sum_{j=0}^{N}X^j h_j(D)
\]
in the normal form of Proposition~\ref{prop:ore}.  Since $F=-D$, the relation $[H,F]=-2F$ is equivalent to $[H,D]=-2D$.  Because $[X^j,D]=-jX^{j-1}$, comparison of normal forms gives
\[
 h_1(D)=2D,
 \qquad
 h_j(D)=0\quad(j\ge2).
\]
Thus $H=2XD+\phi(D)$ for a unique $\phi\in R$. Similarly, writing $E=\sum_{j=0}^{M}X^j e_j(D)$, the relation $[E,F]=H$, or $-[E,D]=H$, forces
\[
 E=X^2D+X\phi(D)+\psi(D)
\]
for a unique $\psi\in R$.  For readability, we record the normal-order calculation.  The identities
\[
 [XD,X]=X,
 \qquad
 [XD,D]=-D,
 \qquad
 [\phi(D),X]=\phi'(D),
 \qquad
 [XD,r(D)]=-Dr'(D)
\]
give
\[
\begin{aligned}
 [2XD,X^2D]&=2X^2D,\\
 [2XD,X\phi]&=2X\phi-2XD\phi',\\
 [2XD,\psi]&=-2D\psi',\\
 [\phi,X^2D]&=2X\phi'D+D\phi'',\\
 [\phi,X\phi]&=\phi\phi'.
\end{aligned}
\]
Since $D$ commutes with $\phi'$, the two terms containing $X D\phi'$ cancel.  Hence
\[
 [H,E]=2X^2D+2X\phi-2D\psi'+D\phi''+\phi\phi'.
\]
Therefore, $[H,E]=2E$ is equivalent to
\begin{equation}\label{eq:psi-ode}
 2D\psi'(D)+2\psi(D)=D\phi''(D)+\phi(D)\phi'(D).
\end{equation}
The integration step is internal to $R$.  Indeed
\[
 (2D\psi)'=2D\psi'+2\psi,
\]
whereas
\[
 \left(D\phi'-\phi+\frac12\phi^2\right)'
 =D\phi''+\phi\phi'.
\]
Thus the two sides of \eqref{eq:psi-ode} are total derivatives.  Because $R\subset\K(D)$ and $\K$ has characteristic zero, the kernel of $\partial_D:R\to R$ is exactly $\K$.  Hence
\[
 2D\psi=D\phi'-\phi+\frac12\phi^2+C
\]
for some $C\in\K$.  Evaluating at $D=0$ gives
\[
 C=\phi(0)-\frac12\phi(0)^2.
\]
Set
\begin{equation}\label{eq:recover-lambda-g}
 \lambda=-\phi(0),
 \qquad
 g(D)=\frac{\phi(D)-\phi(0)}{2D}.
\end{equation}
By Lemma~\ref{lem:evaluation}, the numerator in the definition of $g$ lies in $DR$, so $g\in R$.  Thus
\[
 \phi(D)=2Dg(D)-\lambda.
\]
Substituting this identity into the integrated equation gives, without division by any element outside $R$
\[
 2D\psi=2D^2g'+2D^2g^2-2\lambda Dg,
\]
and hence
\[
 \psi(D)=Dg'(D)+Dg(D)^2-\lambda g(D).
\]
Since $g(D)$ commutes with $D$,
\[
 (X+g)^2D-\lambda(X+g)
 =X^2D+X(2Dg-\lambda)+Dg'+Dg^2-\lambda g=E,
\]
and $H=2(X+g)D-\lambda$.  The formulas \eqref{eq:recover-lambda-g} show uniqueness. Conversely, if $Y=X+g(D)$, then $[D,Y]=1$, so $[YD,Y]=Y$ and $[YD,D]=-D$.  Therefore
\[
 [2YD,Y^2D]=2Y^2D,
 \qquad
 [2YD,-\lambda Y]=-2\lambda Y,
\]
and hence
\[
 [2YD-\lambda,Y^2D-\lambda Y]=2(Y^2D-\lambda Y).
\]
Similarly
\[
 [2YD-\lambda,-D]=-2(-D),
 \qquad
 [Y^2D-\lambda Y,-D]=2YD-\lambda.
\]
Thus \eqref{eq:normal-form} defines an $\mathfrak{sl}_2$-triple.  Finally, the images in \eqref{eq:shear-automorphism} satisfy
\[
 [\sigma_g(D),\sigma_g(X)]=[D,X+g(D)]=1,
\]
so Proposition~\ref{prop:ore} gives an endomorphism $\sigma_g$ of $\algA$.  The composition laws $\sigma_g\sigma_h=\sigma_{g+h}$ and $\sigma_0=\id$ show that its inverse is $\sigma_{-g}$, and applying $\sigma_g$ to the standard triple gives \eqref{eq:normal-form}.
\end{proof}

\end{document}
